\lysh{13171}{4}

\begin{alist}
\item Προφανώς $\alpha_1 = S_1 = 2 \cdot 1^2 + 3 \cdot 1 = 5$.

\item Θέτοντας όπου $\nu$ το $\nu - 1$, παίρνουμε: $S_{\nu-1} = 2(\nu - 1)^2 + 3(\nu - 1) =$
$2(\nu^2 - 2\nu + 1) + 3\nu - 3 = 2\nu^2 - 4\nu + 2 + 3\nu - 3 = 2\nu^2 - \nu - 1$ για κάθε $\nu \ge 2$.

\item Για κάθε $\nu \ge 2$, έχουμε $\alpha_\nu = (\alpha_1 + \alpha_2 + \dots + \alpha_{\nu-1} + \alpha_\nu) - (\alpha_1 + \alpha_2 + \dots + \alpha_{\nu-1}) =$
\[S_\nu - S_{\nu-1} = 2\nu^2 + 3\nu - (2\nu^2 - \nu - 1) = 2\nu^2 + 3\nu - 2\nu^2 + \nu + 1 = 4\nu + 1.\]
Αλλά $\alpha_1 = 5 = 4 \cdot 1 + 1$. Ώστε $\alpha_\nu = 4\nu + 1$, για κάθε $\nu \ge 1$.

\item Για να είναι η ακολουθία $(\alpha_\nu)$ αριθμητική πρόοδος, θα πρέπει η διαφορά
δύο οποιωνδήποτε διαδοχικών όρων να είναι σταθερή.
Πράγματι: $\alpha_{\nu+1} - \alpha_\nu = [4(\nu + 1) + 1] - [4\nu + 1] = 4\nu + 4 + 1 - 4\nu - 1 = 4$.
Άρα η διαφορά $\omega$ είναι ίση με $4$.
\end{alist}

